\section{Comparison and boundary geometry}
\subsection{Adaptive discrimination and description}
Salek--Hayashi--Winter \cite[Sections~II--V]{SHW68} study asymptotic
binary channel-discrimination exponents. Their results include equal
adaptive and nonadaptive exponents for classical--quantum channels,
and separations for some entanglement-breaking channels. Our object
is instead a covering number of a supplied family in the finite-use
metric $d_N$. Its lower packing may use a different binary test for
each pair; it does not supply a universal measurement which identifies
every codeword. The exact reduction in
Lemma~\ref{lem:productreplacer68} uses input erasure, not merely
entanglement breaking, and implies no general absence of an adaptive
advantage for entanglement-breaking channels.

Cooney--Mosonyi--Wilde \cite{CMW70} is a direct antecedent for the
input-erasure mechanism. Its two-replacer discussion and asymptotic
channel-versus-replacer exponents are distinct from an exact finite-use
covering problem. The retraction of Theorem~\ref{thm:retraction70} concerns
admissible centres, not an improvement to those discrimination exponents.
The coherent amplification used here likewise has the established unitary
discrimination antecedents \cite{Acin70,DFY70}.

The two-point programme method has a metrological antecedent in
\cite{DKG67}. Likewise the half-parameter scale has antecedents in
quantum population compression \cite{YBCH68}; there the input consists
of unknown physical copies and the memory interface may include
quantum storage. Here the encoder is given classical matrix data and
the decoder specifies an operation or state. These are different
coding tasks, even when a logarithmic coefficient has the same form.

Mele--Bittel \cite{MB67} and Oufkir--Girardi \cite{OG67} investigate
learning unknown channels in diamond distance. The present encoder
receives its target and asserts no query or sample bound. The
reference-sensitive physical classical-simulation problem of
Naik--Zartab--Gisin--Banik \cite{NZGB67} is different again: a rational
description with a diamond guarantee remains a description of a
quantum operation, not a finite-classical-message implementation
accepting an unknown entangled input. No bound here is asserted for
the hardware dimension of a universal programmable processor.

\subsection{Ranks, tangent directions and remaining joint regimes}
For a general instrument $J$, let $P_y$ project onto the support of
$J_y$. The one-sided tangent cone of the legal instrument body is
\[
 T_{\mathcal C}(J)=\overline{\{t(L-J):t\ge0,\ L\in\mathcal C\}}.
\]
Every element $H$ of this cone satisfies the marginal constraint and
$P_y^\perp H_yP_y^\perp\ge0$. These necessary conditions identify
where positivity becomes active; by themselves they assert no reuse
modulus. A boundary classification must examine the metric of
repeated use along support-changing directions and along reversible
coherent directions, rather than count affine coordinates alone.
Proposition~\ref{prop:boundaryphase67} shows that a coherent unitary
phase can have scale $N^{-1}$, whereas the preparation family in
Theorem~\ref{thm:rankentropy68} has an $N^{-1/2}$ coding scale even
at rank-deficient outputs. In that preparation family the Choi ranks
are $d\,\operatorname{rank}\sigma_y$, and the rank bounds in the
theorem are output ranks. They are not arbitrary Choi-rank strata of
all instruments.

Theorem~\ref{thm:coherententropy70} now gives a joint law in the
non-input-erasing tensor family $UXU^*\otimes\sigma_y$. It includes a
full projective unitary manifold and support-changing preparation blocks,
but no input-dependent measurement probabilities. Its mixed coefficient
$u+v/2$ is established by two explicit adaptive moduli and matching product
packings, not by a tangent-dimension count alone. These qualifications are
essential: it is not a classification of all coherent/noisy couplings.

Theorems~\ref{thm:jointcode68} and~\ref{thm:rankentropy68} settle the
joint $N,\delta$ dependence in their specified fixed-dimensional
families. They do not assert dimension-uniform constants, a uniform
limit as the positive margin tends to zero, or a full tangent-cone
classification for general disturbing instruments. The older
exponential-accuracy width result retains its multiplicative gap in
the separate label-width resource. None of these description laws
is a mutable-workspace converse for general variable-input
trajectory simulation.
